\section{Introduction}
\label{sec:intro}

Hematoxylin-and-eosin (H\&E) staining of tissue sections is the most widely used
assay in cancer medicine~\cite{bera2019}. Every diagnostic
slide records tissue composition, cell densities, nuclear morphology and the
spatial arrangement of tumor and microenvironmental
cells~\cite{gurcan2009, vanderlaak2021}. For genomics, transcriptomics and
proteomics, pan-cancer resources release standardized, queryable
tables~\cite{kandoth2013, hoadley2014, li2023cptac, corces2018}; no equivalent
resource exists for histopathology. Diagnostic slides are archived at scale, but
the quantities a pathologist assesses from them are rarely recorded as numerical
data.

Bulk molecular assays average over the whole sample and do not record the
spatial organization of the tissue. Yet the density of immune cells in the tumor
core and at the invasive margin is prognostic in several cancer
types~\cite{galon2006, fridman2012immune, pages2018}. Saltz et~al.\ showed that
deep learning applied to public pathology slides recovers tumor-infiltrating
lymphocyte density at pan-cancer scale~\cite{saltz2018til}, but a density
computed over the whole slide does not distinguish lymphocytes inside tumor
epithelium from lymphocytes confined to the surrounding stroma, the pattern that
defines immune exclusion. Existing public resources each cover part of the
problem: cBioPortal provides molecular and clinical tables without quantitative
morphology~\cite{cerami2012cbio}, the Cancer Imaging Archive provides images
without associated statistical analyses~\cite{clark2013tcia}, and the Human
Protein Atlas provides immunohistochemistry-based maps of individual
proteins~\cite{uhlen2015hpa, uhlen2017pathology}, which are not available for
routine H\&E slides.

Computational pathology has made the extraction of tissue features tractable.
Early morphometric studies linked handcrafted image features to outcome and
molecular alterations in single cancer types~\cite{beck2011, cooper2015, yu2016}; more recent systems
predict molecular alterations, expression programs and outcomes directly from
H\&E~\cite{coudray2018, schmauch2020he2rna, bannier2024}, as reviewed
in~\cite{echle2021}; and foundation models now produce tile-level embeddings
that can be aggregated to the slide
level~\cite{chen2024uni, vorontsov2024virchow, hoptimus0}. These results
show that H\&E images contain substantial biological information. Foundation-model
embeddings are information-dense representations that support many downstream
tasks, including mutation, biomarker and outcome
prediction~\cite{filiot2023phikon, filiot2024phikonv2, bannier2024, bannier2024her2},
but they are not interpretable: an individual coordinate does not
correspond to a defined tissue property, so an association found with an
embedding cannot be attributed to specific cells or tissue regions.
Interpretable feature sets have been proposed~\cite{diao2021hif, abel2024}, but
no pan-cancer set of interpretable measurements has been released together with
its molecular and clinical associations and evaluated in an independent cohort.

Image statistics depend on staining protocol, scanner, fixation and case
selection, so a measurement defined in one cohort can have a different
distribution, and a different relationship to molecular variables, in another.
Measurements and associations derived from The Cancer Genome Atlas (TCGA)
therefore need to be re-estimated on slides from other institutions. The Clinical
Proteomic Tumor Analysis Consortium (CPTAC) provides independent diagnostic slides
with matched transcriptomic and proteomic profiles across several cancer
types~\cite{li2023cptac}. This allows two tests: whether the measurement
distributions and their correlations with gene expression are reproduced in an
independent cohort, and whether the same correlations are observed at the
protein level.

Here we present HistoAtlas. From \num{6745} TCGA diagnostic slides from
\num{6088} patients spanning 21 solid-tumor types, we derive 38 interpretable,
compartment-resolved histomic measurements and compute their associations with
survival, gene expression, copy number, pathway activity and immune state,
reporting effect sizes, confidence intervals, minimum detectable effects and
explicitly declared multiple-testing families throughout. We show that the 38
measurements can be summarized by six continuous programs, each correlated with
specific MSigDB Hallmark pathways~\cite{liberzon2015}; that the measurement distributions, their
correlations with gene expression and two of the six programs are reproduced in
\num{1095} independent CPTAC slides; that the correlations with gene expression
are also observed at the protein level, with 132 measurement--protein
correlations remaining significant after adjustment for the matched transcript;
and that adding the measurements to survival models that already include
clinical covariates and matched transcripts slightly improves discrimination. The resource is released as an interactive atlas and
a programmatic interface, in which every reported statistic is linked to the
tissue compartments and individual cells from which it was computed.
